\documentclass[../../main]{subfiles}

\ifSubfilesClassLoaded{%
    \zexternaldocument*{../1.Category-Theory/exported/aux/main}%
    \zexternaldocument*{../2.Homological-Algebra/exported/aux/main}%
    \zexternaldocument*{../3.Sheaves/exported/aux/main}%
    \zexternaldocument*{1.Cohomology-of-sheaves}%
    \zexternaldocument*{2.Cohomology-and-homotopy}%
}{}

\begin{document}
\chapter{Some examples of cohomology}

We focus our attention now to the computation of the cohomology of some topological spaces.
Our main tools will be the \hyperref[thm:mayer-vietoris-open]{Mayer-Vietoris sequence} and the \hyperref[thm:homotopy-invariance-cohomology]{homotopy invariance of cohomology}.

\begin{example}[\(\mathbb{S}^{1}\)]\label{ex:cohomology-S1}
  Let \(\sheaf{F} \in \Mod( \underline{\mathbf{k}}_{S} )\) be a \hyperref[def:sheaf]{sheaf of \(\mathbf{k}\)-modules} over \(S\) and \(X \coloneq S \times \mathbb{S}^{1}\).
  
  Consider the open cover of \(\mathbb{S}^{1} = H^{+} \cup H^{-}\) where \(H^{+}\) is the upper hemisphere and \(H^{-}\) the lower one, both extended slightly past the equator.
  Let \(U = S \times H^{+}\) and \(V = S \times H^{-}\).
  Since \(H^{+}\) and \(H^{-}\) are homeomorphic to open intervals, which are contractible, both \(U\) and \(V\) are \hyperref[def:homotopy-over]{contractible over \(S\)}.
  Furthermore \(U \cap V = \left( S \times W_{1} \right) \coprod \left( S \times W_{2} \right)\) consists of two disjoint pieces, each of these \hyperref[def:homotopy-over]{contractible over \(S\)}.

  Therefore since \hyperref[cor:cohomology-and-contractible]{cohomology of contractible spaces is concentrated in degree 0} we only have
  \begin{align*}
    \operatorname{R}^{0} p_{U, \ast}\sheaf{F}|_{U} &\cong \sheaf{F} \\
    \operatorname{R}^{0}p_{V, \ast}\sheaf{F}|_{V} &\cong \sheaf{F} \\
    \operatorname{R}^{0}p_{U \cap V, \ast} \sheaf{F}|_{U \cap V} &\cong \sheaf{F} \oplus \sheaf{F}
  \end{align*}
  Writing the \hyperref[thm:mayer-vietories-open]{Mayer-Vietoris sequence} since the higher terms vanish we end up with
  \begin{center}
    \begin{tikzcd}
      0 \arrow[r] & \operatorname{R}^{0} p_{X, \ast} \sheaf{F}_{X} \arrow[r] & \sheaf{F} \oplus \sheaf{F} \arrow[r, "\alpha"] & \sheaf{F} \oplus \sheaf{F} \arrow[r, "\delta"] & \operatorname{R}^{1} p_{X, \ast} \sheaf{F}|_{X} \arrow[r] & 0
    \end{tikzcd}
  \end{center}
  Therefore directly we get \(\operatorname{R}^{i}p_{X, \ast}\sheaf{F}_{X} = 0\) for all \(i \geq 2\).

  Note that \(\alpha\colon \sheaf{F} \oplus \sheaf{F} \to \sheaf{F} \oplus \sheaf{F}\) maps a pair of sections \((s_{U}, s_{V})\) into their difference on the overlap.
  Since the section \(s_{U}\) is defined globally on \(U\) its restriction to \(W_{1}\) and to \(W_{2}\) is identical and analogously for \(s_{V}\).
  Therefore we have
  \begin{equation*}
    \alpha(s_{U}, s_{V}) = (s_{U} - s_{V}, s_{U} - s_{V})
  \end{equation*}
  and thus \(\ker \alpha = \{(s, s) \mid s \in \sheaf{F}\}\) is the diagonal which means that
  \begin{equation*}
    \operatorname{R}^{0}p_{X, \ast} \sheaf{F}_{X} \cong \sheaf{F}
  \end{equation*}

  On the other hand by \hyperref[def:exact-sequence]{exactness} \(\delta\) is surjective, meaning that
  \begin{equation*}
    \operatorname{R}^{1}p_{X, \ast}\sheaf{F}|_{X} \cong \coker \alpha \cong \quot{\sheaf{F} \oplus \sheaf{F}}{\operatorname{Im} \alpha} = \quot{\sheaf{F} \oplus \sheaf{F}}{\langle(s, s)\rangle} \cong \sheaf{F}
  \end{equation*}

  Therefore
  \begin{equation*}
    \operatorname{R}^{i}p_{S \times \mathbb{S}^{1}, \ast}\sheaf{F}_{X} = \begin{cases}
      \sheaf{F} & i = 0 \\
      \sheaf{F} & i = 1 \\
      0 & i \geq 2
    \end{cases}
  \end{equation*}
  In particular taking \(S = \{\ast\}\) since \(\sheaf{F} \cong \underline{M}\) for some \(M \in \Mod(\mathbf{k})\) we get
  \begin{equation*}
    \cohomology{i}(\mathbb{S}^{1}; \underline{M}) = \begin{cases}
      M & i = 0 \\
      M & i = 1 \\
      0 & i \geq 2
    \end{cases}
  \end{equation*}
\end{example}

\begin{example}[\(\mathbb{S}^{n}\)]
  Let \(\sheaf{F} \in \Mod(\underline{\mathbf{k}}_{S})\).
  We are going to see that cohomology of \(X = S \times \mathbb{S}^{n}\) over \(S\) by induction over \(n\) is
  \begin{equation*}
    \operatorname{R}^{i}p_{X, \ast}\sheaf{F}|_{X} = \begin{cases}
      \sheaf{F} & i = 0 \\
      \sheaf{F} & i = n \\
      \sheaf{F} & i \notin \{0, n\}
    \end{cases}
  \end{equation*}
  We have already computed the base case \(n = 1\) in \autoref{ex:cohomology-S1}, so we can suppose \(n \geq 2\).
  
  Consider the open cover \(\mathbb{S}^{n} = H^{+} \cup H^{-}\) where \(H^{+}\) is the open upper hemisphere and \(H^{-}\) is the lower open hemisphere, both extending slightly past the equator.
  We have that \(U \coloneq S \times H^{+}\) and \(V \coloneq S \times H^{-}\) are \hyperref[def:homotopy-over]{contractible over \(S\)}.
  Furthermore since \(H^{+} \cap H^{-}\) is a thickened equatorial band it can be \textit{squeezed} into the equator itself, which it is a \((n-1)\)-sphere.
  This means that \(H^{+} \cap H^{-}\) \hyperref[def:deformation-retract]{strong-deformation retracts} onto \(\mathbb{S}^{n-1}\) and thus, \(U \cap V\) \hyperref[def:deformation-retract]{strong-deformation retracts} onto \(S \times \mathbb{S}^{n-1}\).

  Since \(\mathbb{S}^{n-1}\) is connected the first terms of the \hyperref[thm:mayer-vietoris-open]{Mayer-Vietoris sequence} are
  \begin{center}
    \begin{tikzcd}
      0 \arrow[r] & \operatorname{R}^{0} p_{X, \ast} \sheaf{F}_{X} \arrow[r] & \sheaf{F} \oplus \sheaf{F} \arrow[r, "\alpha"] & \sheaf{F} \arrow[r, "\delta"] & \operatorname{R}^{1} p_{X, \ast} \sheaf{F}|_{X} \arrow[r] & 0 \oplus 0 \arrow[r] & \cdots
    \end{tikzcd}
  \end{center}
  In this case we have \(\alpha(s_{U}, s_{V}) = s_{U} - s_{V}\) and thus \(\operatorname{R}^{0}p_{X, \ast}\sheaf{F}|_{X} \cong \sheaf{F}\).
  Furthermore since \(\alpha\) is clearly surjective and \(\delta\) is surjective because the \hyperref[def:exact-sequence]{sequence is exact} we have \(\operatorname{R}^{1}p_{X, \ast}\sheaf{F}|_{X} \cong 0\).

  If we focus now on the higher terms of the \hyperref[thm:mayer-vietoris-open]{Mayer-Vietoris sequence} we have, for \(i \geq 1\)
  \begin{center}
    \begin{tikzcd}
      \cdots \arrow[r] & \operatorname{R}^{i}p_{U, \ast} \oplus \operatorname{R}^{i}p_{V, \ast} \arrow[r] & \operatorname{R}^{i}p_{U \cap V, \ast} \arrow[r, "\delta"] & \operatorname{R}^{i+1}p_{X, \ast} \arrow[r] & \operatorname{R}^{i+1}p_{U, \ast} \oplus \operatorname{R}^{i+1}p_{V, \ast} \arrow[r] & \cdots
    \end{tikzcd}
  \end{center}
  Since \(U\) and \(V\) are \hyperref[def:homotopy-over]{contractible over \(S\)} these terms are zero by \autoref{cor:cohomology-and-contractible} and thus we get an \hyperref[def:categorical-isomorphism]{isomorphism}
  \begin{equation*}
    \operatorname{R}^{i}p_{U \cap V, \ast} \isoarrow \operatorname{R}^{i+1}p_{X, \ast}
  \end{equation*}
  Substituting \(U \cap V\) we get
  \begin{equation*}
    \operatorname{R}^{i+1}p_{X, \ast} \sheaf{F}|_{S \times \mathbb{S}^{n}} \cong \operatorname{R}^{i}p_{X, \ast} \sheaf{F}|_{S \times \mathbb{S}^{n-1}} \cong \begin{cases}
      \sheaf{F} & i \in \{0, n-1\} \\
      \sheaf{F} & i \notin \{0, n-1\}
    \end{cases}
  \end{equation*}
  
  Therefore
  \begin{equation*}
    \operatorname{R}^{i}p_{S \times \mathbb{S}^{n}, \ast}\sheaf{F}|_{S \times \mathbb{S}^{n}} = \begin{cases}
      \sheaf{F} & i = 0 \\
      \sheaf{F} & i = n \\
      0 & i \notin \{0, n\}
    \end{cases}
  \end{equation*}
  In particular taking \(S = \{\ast\}\) since \(\sheaf{F} \cong \underline{M}\) for some \(M \in \Mod(\mathbf{k})\) we get
  \begin{equation*}
    \cohomology{i}(\mathbb{S}^{n}; \underline{M}) = \begin{cases}
      M & i = 0 \\
      M & i = n \\
      0 & i \notin \{0, n\}
    \end{cases}
  \end{equation*}
\end{example}

\begin{example}[\(\mathbb{R}^{n} \setminus \{0\}\)]
  We are going to see first that \(X \coloneq S \times \left( \mathbb{R}^{n} \setminus \{0\} \right)\) is a \hyperref[def:deformation-retract]{strong deformation retract} of \(A \coloneq S \times \mathbb{S}^{n-1}\).
  To do this, consider the inclussion \(i\colon A \hookrightarrow X\) and the retraction \(r\colon X \to A, r(s, v) \coloneq \left( s, \frac{v}{\lVert v\rVert} \right)\).
  As the \hyperref[def:homotopy-over]{homotopy} we can take
  \begin{align*}
    H\colon& X \times [0, 1] \to A \\
    &(s, v, t) \mapsto \left( s, (1 - t)v + t \frac{v}{\lVert v \rVert} \right)
  \end{align*}
  Therefore by \autoref{cor:cohomology-and-retraction} for all \(\sheaf{F} \in \Mod(\underline{\mathbf{k}}_{S})\) we have and \(i \geq 0\)
  \begin{equation*}
    \operatorname{R}^{i}p_{X, \ast}\sheaf{F}|_{X} \cong \operatorname{R}^{i}p_{A, \ast}\sheaf{F}|_{X}
  \end{equation*}
  In particular if \(n \geq 2\) we have
  \begin{equation*}
    \operatorname{R}^{i}p_{S \times \left(\mathbb{R}^{n} \setminus \{0\}\right), \ast}\sheaf{F}|_{S \times \left( \mathbb{R}^{n} \setminus \{0\} \right)} \cong \begin{cases}
      \sheaf{F} & i = 0 \\
      \sheaf{F} & i = n-1 \\
      0 & i \notin \{0, n-1\}
    \end{cases}
  \end{equation*}
  The case \(n = 1\) is different since \(\mathbb{R} \setminus \{0\}\) has two connected components.
  We have a \hyperref[def:deformation-retract]{deformation retract} \(\mathbb{R} \setminus \{0\} \to \mathbb{S}^{0} \cong \{\pm 1\}\).
  Therefore
  \begin{equation*}
    \operatorname{R}^{i}p_{S \times \left( \mathbb{R} \setminus \{0\} \right), \ast}\sheaf{F}|_{S \times \left( \mathbb{R} \setminus 0 \right)} \cong \begin{cases}
      \sheaf{F} \oplus \sheaf{F} & i = 0 \\
      0 & i \neq 0
    \end{cases}
  \end{equation*}

  In particular taking \(S = \{\ast\}\) since \(\sheaf{F} \cong \underline{M}\) for some \(M \in \Mod(\mathbf{k})\) we get
  \begin{align*}
    \cohomology{i}(\mathbb{R}^{n} \setminus \{0\}; \underline{M}) &= \begin{cases}
      M & i = 0 \\
      M & i = n-1 \\
      0 & i \notin \{0, n-1\}
    \end{cases} \\
     \cohomology{i}(\mathbb{R} \setminus \{0\}; \underline{M}) &= \begin{cases}
      M \oplus M & i = 0 \\
      0 & i \neq 0
    \end{cases}
  \end{align*}
\end{example}

\end{document}
